\begin{textsample}{NEG Dim 6 – global\_north\_2024\_12 – Score -1.00 – global\_north\_2024\_12/118955846.txt}  \label{ex:f6_neg_196}
Editorial : Diplomacy ruptured by regrettable Israeli embassy closure Israeli ' s embassy to Ireland.
Photo : Arthur Carron Editorial Mon 16 Dec 2024 at 03:30 Three-quarters of a century ago, in a globe that was recovering from World War II and facing into the Cold War, the then US secretary of state Dean Acheson summed up why it is crucial to remain diplomatically engaged, even in hostile or dangerous parts of the world."We maintain diplomatic relations with other countries primarily because we are all on the same planet and must do business with each other.
We do not establish an embassy in a foreign? country to show approval of its government.
We do so to have a channel through which to conduct essential? government relations and to protect legitimate United States interests,"Acheson wrote.
Seventy-five years later, Israel is not heeding that message.
There is no clear statistical measure of the value to a country of having a diplomatic presence in another nation.
Certainly it makes it easier for citizens of that country to live, work and integrate in the adopted country if there is such a presence.
Arguably, there is a trade advantage to be able to make introductions between producers and consumers on either side of the enterprise chain and learn how to reduce the red tape at either end.
Moreover, it shows a desire to develop a positive relationship in good times -- or bad.
The bulk of these interactions between an embassy and the host country are civil, polite and easy -- diplomatic, by their nature.
But when difficult words need to be expressed, it allows such communications to be clearly understood.
Ireland ' s stance on Israel ' s actions in Gaza is not driven by malice The traditionally strong ties between Ireland and?
Israel have become strained over the last 15 months as a result of the war in Gaza.
The Government was among the strongest voices to condemn the Hamas attacks on innocent Israeli citizens on October 7, with the bloodshed, murders and abductions shocking many in this country.
The Government has also been clear that Israel has a right to defend itself but that the response had to be proportionate.
What has been witnessed in Gaza over the last 15 months has not been proportionate, with unprecedented death, widespread destruction of property and restriction of humanitarian aid.
The Government has chosen to recognise the Palestinian state but not to expel the Israeli ambassador, as some opposition parties have demanded.
Israel has now decided to rupture diplomatic relations by closing down its embassy in Dublin.
The Israeli government cites what it calls"the extreme anti-Israel policies of the Irish Government".
The inciting incident appears to be the announcement of support for South Africa ' s legal action against Israel in the International Court of Justice, accusing Israel of"genocide".
The \textbf{caretaker} Government has described the move as a"deeply regrettable decision".
Rather than making all sides take stock and seek to set their differences aside, the danger is that an embassy closure can also lead to more entrenched policies of diplomatic disengagement.
Ireland ' s stance on Israel ' s actions in Gaza is not driven by malice.
Sometimes the best friends are those who are not afraid to tell uncomfortable truths.

% matched lemmas: caretaker
\end{textsample}
